\section{ChartRWC 数据集构建}

本章介绍 ChartRWC 数据集的需求、构建方法以及数据集统计信息。

\subsection{数据集需求}

现有数据驱动式图表推荐方法的成功，在很大程度上依赖于数据集质量~\cite{dibia2019data2vis,hu2019vizml}。本文对学术界常用图表数据集进行了深入调研，例如 ChartQA~\cite{masry2022chartqa}、ChartBench~\cite{Xu2023ChartBench} 和 ChartGen~\cite{kondic2025chartgen}，并识别出以下三个局限：

\begin{itemize}
    \item \textbf{缺乏明确的用户意图：} 这些数据集主要由“数据-图表”对构成，但几乎完全缺少关于选择特定图表背后分析意图的标注。因此，模型难以学习数据与用户深层分析目标之间的映射关系。
    \item \textbf{图表类型分布严重不均：} 在真实世界的数据可视化中，柱状图、折线图和散点图占据了绝大多数实例。因此，现有数据集往往呈现极端类别不平衡，使模型或检索方法在推荐箱线图、雷达图、漏斗图等较少使用的图表类型时表现不佳。
    \item \textbf{数据来源相对同质：} 许多数据集中的表格主要来源于商业报表或科学出版物等特定领域，这限制了模型在处理更广泛领域数据时的泛化能力，例如社交媒体、个人健康等场景。
\end{itemize}

为克服上述不足，本文为 ChartRWC 数据集设定如下三个数据集需求：

\ding{172} \textbf{DR1：足够体量。} 数据集应包含足够的样本规模，以满足检索方法与场景覆盖对数据量的需求。

\ding{173} \textbf{DR2：多样性与平衡性。} 数据集需覆盖真实世界中常见图表类型，并使不同图表类型之间的样本体量保持相对平衡。

\ding{174} \textbf{DR3：广泛的领域覆盖。} 数据集应具有广泛的领域覆盖范围，以增强其泛化能力。

\subsection{数据集构建方法}

为达成第4.1节设定的设计目标，本文采用包含“种子数据收集”与“增强与平衡”的两阶段方法来构建 ChartRWC 数据集。

\subsubsection{种子数据收集}

种子数据主要来源于现有高质量学术数据集，以确保初始数据的准确性和规范性。本文使用公开发布的数据集 ChartGen~\cite{kondic2025chartgen}。该完整数据集包含18种图表类型，并包含结构化数据、图表图像、图表摘要、DocTags、问答对以及代码。

为构建一个能够反映真实世界数据可视化实践并作为稳健基准的数据集，本文首先需要定义本研究的核心图表类型范围。本文没有直接采用 ChartGen 中的全部图表类型，因为其内部样本分布存在显著采集偏差，并且包含许多定义重叠或过于专业化的变体。本文从该数据集中筛选了10种图表类型，并提取源数据、图表图像、问答对和图表描述等核心元素。通过这一流程，最终构建了一个由154,411个样本组成的种子数据集。尽管通过筛选高质量源数据集的方式为种子数据奠定了良好基础，但它仍存在第4.1节所提到的两个关键缺陷。

\textbf{图表类别平衡性不足。}
如表\ref{tab:Statistic} 第二列和第三列所示，最常见的四种图表类型，即折线图、柱状图、箱线图和散点图，占全部种子数据的74.3\%，形成了显著的长尾分布。若以此作为基准，检索方法的整体性能指标可能会被这些高频类别主导。这不仅无法准确评估稀有图表类型上的检索能力，也会削弱其作为公平、全面评测基准的有效性。

\textbf{文本描述语义贫乏。}
绝大多数描述遵循固定模板，其功能仅限于简单识别图表组成元素，而未能传达图表所揭示的核心洞见或叙事。这种语义贫乏会对两种主流检索范式构成明显障碍。对于稀疏检索而言，描述性词汇不足会导致召回率较低；对于稠密检索而言，生成的嵌入向量缺乏区分度，无法与包含真实用户意图的查询在语义空间中有效匹配。

\begin{table}[htbp]
\centering
\small
\resizebox{\textwidth}{!}{
\begin{tabular}{lcccc}
\hline
图表类型 & ChartGen数量 & ChartGen占比(\%) & ChartRWC数量 & ChartRWC占比(\%) \\ \hline
Bar Chart & 28,934 & 18.74 & 17,494 & 9.97 $\downarrow$ \\
Box Plot & 24,664 & 15.97 & 17,498 & 9.98 $\downarrow$ \\
Bubble Plot & 12,188 & 7.89 & 17,482 & 9.97 $\uparrow$ \\
Funnel Chart & 3,657 & 2.37 & 17,293 & 9.86 $\uparrow$ \\
Line Chart & 42,218 & 27.34 & 17,495 & 9.97 $\downarrow$ \\
Pie Chart & 12,969 & 8.40 & 17,486 & 9.97 $\uparrow$ \\
Radar Chart & 5,185 & 3.36 & 17,468 & 9.96 $\uparrow$ \\
Scatter Plot & 18,837 & 12.20 & 17,497 & 9.97 $\downarrow$ \\
Stacked Bar Chart & 1,269 & 0.82 & 16,242 & 9.26 $\uparrow$ \\
Treemap Chart & 4,490 & 2.91 & 19,453 & 11.09 $\uparrow$ \\ \hline
\end{tabular}
}
\caption{ChartGen 与 ChartRWC 的统计信息}
\label{tab:Statistic}
\end{table}

\subsubsection{增强与平衡}

由于种子数据集在图表类型样本数量平衡性和图表文本描述语义丰富度方面存在核心缺陷，本文设计并执行了一套策略性数据集扩展与平衡流程，以达成第4.1节中设定的三项设计目标。该流程包含三种策略：

\ding{172} \textbf{增强数据集异构性。}
为进一步丰富数据的内在多样性并满足异构性目标，本文对种子数据表本身进行了变换，包括列子集选择（通过选取原表列子集创建新表）和数据类型转换（例如，将连续数值变量分箱为类别变量），以增加数据表数量。

\ding{173} \textbf{增加文本描述多样性。}
为增加语义多样性，本文构建了一个主题表，并从中随机抽取主题。随后，将这些主题与图表数据结合，利用 GPT-4o 生成自然语言描述~\cite{chatgpt4o2024}。这些新增描述覆盖了多个真实世界数据分析应用场景。

\ding{174} \textbf{平衡图表类型体量。}
为解决种子数据中的不平衡问题，本文通过减少占比较高图表类型的数量，并增加稀有图表类型的样本数量，缓解数据集中的长尾分布问题。本文设定每种图表类型样本数量约为17,000。对于种子数据中过度表示的图表类型，采用无放回随机下采样方法将其削减至目标比例；对于样本数量不足的图表类型，则使用前述策略\ding{172}和策略\ding{173}生成样本。

该合成过程借助 Matplotlib~\cite{Hunter2007Matplotlib} 和 D3~\cite{bostock2011d3} 等可视化库，基于预定义模板和结构化数据渲染有效样本。通过上述流程，数据集平衡性得到了改善。如表\ref{tab:Statistic} 第三列和第五列所示，占比最低的图表类型从0.82\%提升至9.26\%，而占比最高的图表类型从27.34\%降低至9.97\%。

\subsection{数据集统计与验证}

经过上述两阶段构建流程后，最终 ChartRWC 数据集共包含175,408个样本，如表\ref{tab:Statistic} 所示。其中，部分样本来自种子数据，其余样本由增强与平衡策略生成。每个条目均包含数据表、PNG格式图像、文本描述、问答对以及图表类型标签。

在规模方面，ChartRWC 数据集包含超过17万个样本，足以支撑多模态大语言模型训练，同时也能满足稀疏检索和稠密检索方法的需求（DR1）。在多样性与平衡性方面，如表\ref{tab:Statistic} 所示，与种子数据集相比，ChartRWC 的分布得到了显著改善（DR2）。同时，ChartRWC 数据集也保留并扩展了种子数据的异构性（DR3）。

为验证数据质量，本文邀请2位数据可视化领域参与者（E1、P5）以及3位数据应用领域参与者（P6、P7、P9），对随机抽样的500个样本进行人工评估，每位参与者评估100个样本。评估结果显示，专家认为96.2\%的样本是“合理的”或“非常合理的”，证明了增强与平衡策略的有效性

\newpage
